% GENERATED FILE. Edit canonical lessons, then run npm run generate:books.
% canonical-source-hash: fnv1a64:4412bc07567d3cec
% canonical-lessons: SA-C417-pulinam, SA-C417-phenah, SA-C417-budbudah, SA-C417-sthalam, SA-C417-vela, SA-R417-review-410-417

\chapter{Sandbanks and Shores}
\label{ch:sa-sandbanks-and-shores}
\hlchaptermodality{417}

\begin{chapteropening}
\textbf{By the end of this chapter:} \emph{I can describe the sand, foam and shore by the water.}

Complete the last lesson of chapter 417: I can describe the sand, foam and shore by the water.
\end{chapteropening}

\section[\texorpdfstring{\sk{पुलिनम्} (pulinam)}{पुलिनम् (pulinam)}]{\sk{पुलिनम्} (pulinam) — a sandbank}

\label{lesson:SA-C417-pulinam}



\begin{quote}
Before the new one: say the Sanskrit for rows, then the Sanskrit for comes ashore.
\end{quote}

\subsection*{You'll want to know: \sk{पुलिनम्}}
\textbf{\sk{पुलिनम्}} — \emph{pulinam} — \textquotedblleft{}a sandbank\textquotedblright{}.

\begin{quote}
\sk{नद्याः} \sk{मध्ये} \sk{पुलिनम्} \sk{अस्ति।}
\end{quote}

\begin{quote}
\emph{nadyāḥ madhye pulinam asti.}
\end{quote}

\begin{quote}
There is a sandbank in the middle of the river.
\end{quote}

A neuter noun.

\begin{cousinweb}
\textbf{\sk{पुलिनम्}} (\emph{pulinam}) is the word MW gives for a sandbank, a small island left by a river; MW gives no derivation, so its origin is uncertain (MW).
\end{cousinweb}

\begin{grammarlens}[title={What you've built}]
One more word for getting around.
\end{grammarlens}

\subsection*{Your turn}
\begin{itemize}
\raggedright
  \item \emph{Say it:} \emph{pulinam}
  \item \emph{Say it:} \emph{pulinam}, once more
  \item \emph{Say it:} \emph{nadyāḥ madhye pulinam asti.} — There is a sandbank in the middle of the river.
  \item \emph{Recall:} say the Sanskrit for rows, then the Sanskrit for comes ashore, then say \emph{pulinam} again
\end{itemize}

\begin{tcolorbox}[breakable,colback=teal!4,colframe=teal!35!black,title={Before you move on}]
What does \sk{पुलिनम्} mean? (\textquotedblleft{}A sandbank\textquotedblright{}.) Say it once more.
\end{tcolorbox}
\section[\texorpdfstring{\sk{फेनः} (phenaḥ)}{फेनः (phenaḥ)}]{\sk{फेनः} (phenaḥ) — foam}

\label{lesson:SA-C417-phenah}



\begin{quote}
Before the new one: say the Sanskrit for comes ashore, then the Sanskrit for a sandbank.
\end{quote}

\subsection*{You'll want to know: \sk{फेनः}}
\textbf{\sk{फेनः}} — \emph{phenaḥ} — \textquotedblleft{}foam\textquotedblright{}.

\begin{quote}
\sk{तरंगेषु} \sk{श्वेतः} \sk{फेनः} \sk{अस्ति।}
\end{quote}

\begin{quote}
\emph{taraṅgeṣu śvetaḥ phenaḥ asti.}
\end{quote}

\begin{quote}
There is white foam on the waves.
\end{quote}

A masculine noun.

\begin{cousinweb}
\textbf{\sk{फेनः}} (\emph{phenaḥ}) is the word MW gives for foam, froth; MW connects it, probably, with a root √phaṇ, so its origin is uncertain; MW compares English \emph{foam} and Slavonic \emph{pěna} (MW).
\end{cousinweb}

\begin{grammarlens}[title={What you've built}]
One more word for getting around.
\end{grammarlens}

\subsection*{Your turn}
\begin{itemize}
\raggedright
  \item \emph{Say it:} \emph{phenaḥ}
  \item \emph{Say it:} \emph{phenaḥ}, once more
  \item \emph{Say it:} \emph{taraṅgeṣu śvetaḥ phenaḥ asti.} — There is white foam on the waves.
  \item \emph{Recall:} say the Sanskrit for comes ashore, then the Sanskrit for a sandbank, then say \emph{phenaḥ} again
\end{itemize}

\begin{tcolorbox}[breakable,colback=teal!4,colframe=teal!35!black,title={Before you move on}]
What does \sk{फेनः} mean? (\textquotedblleft{}Foam\textquotedblright{}.) Say it once more.
\end{tcolorbox}
\section[\texorpdfstring{\sk{बुद्बुदः} (budbudaḥ)}{बुद्बुदः (budbudaḥ)}]{\sk{बुद्बुदः} (budbudaḥ) — a bubble}

\label{lesson:SA-C417-budbudah}



\begin{quote}
Before the new one: say the Sanskrit for a sandbank, then the Sanskrit for foam.
\end{quote}

\subsection*{You'll want to know: \sk{बुद्बुदः}}
\textbf{\sk{बुद्बुदः}} — \emph{budbudaḥ} — \textquotedblleft{}a bubble\textquotedblright{}.

\begin{quote}
\sk{जले} \sk{बुद्बुदाः} \sk{क्षणं} \sk{तिष्ठन्ति।}
\end{quote}

\begin{quote}
\emph{jale budbudāḥ kṣaṇaṃ tiṣṭhanti.}
\end{quote}

\begin{quote}
Bubbles last a moment on the water.
\end{quote}

A masculine noun.

\begin{cousinweb}
\textbf{\sk{बुद्बुदः}} (\emph{budbudaḥ}) is the word MW gives for a bubble; MW sets it beside \emph{buḍabuḍa}, the sound of bubbling, and gives no derivation, so its origin is uncertain (MW).
\end{cousinweb}

\begin{grammarlens}[title={What you've built}]
One more word for getting around.
\end{grammarlens}

\subsection*{Your turn}
\begin{itemize}
\raggedright
  \item \emph{Say it:} \emph{budbudaḥ}
  \item \emph{Say it:} \emph{budbudaḥ}, once more
  \item \emph{Say it:} \emph{jale budbudāḥ kṣaṇaṃ tiṣṭhanti.} — Bubbles last a moment on the water.
  \item \emph{Recall:} say the Sanskrit for a sandbank, then the Sanskrit for foam, then say \emph{budbudaḥ} again
\end{itemize}

\begin{tcolorbox}[breakable,colback=teal!4,colframe=teal!35!black,title={Before you move on}]
What does \sk{बुद्बुदः} mean? (\textquotedblleft{}A bubble\textquotedblright{}.) Say it once more.
\end{tcolorbox}
\section[\texorpdfstring{\sk{स्थलम्} (sthalam)}{स्थलम् (sthalam)}]{\sk{स्थलम्} (sthalam) — dry land}

\label{lesson:SA-C417-sthalam}



\begin{quote}
Before the new one: say the Sanskrit for foam, then the Sanskrit for a bubble.
\end{quote}

\subsection*{You'll want to know: \sk{स्थलम्}}
\textbf{\sk{स्थलम्}} — \emph{sthalam} — \textquotedblleft{}dry land\textquotedblright{}.

\begin{quote}
\sk{नौका} \sk{स्थलं} \sk{प्राप्नोति।}
\end{quote}

\begin{quote}
\emph{naukā sthalaṃ prāpnoti.}
\end{quote}

\begin{quote}
The boat reaches dry land.
\end{quote}

A neuter noun.

\begin{cousinweb}
\textbf{\sk{स्थलम्}} (\emph{sthalam}) is the word MW gives for dry land, firm earth as against water, and also a place or spot; MW gives no derivation, so its origin is uncertain (MW).
\end{cousinweb}

\begin{grammarlens}[title={What you've built}]
One more word for getting around.
\end{grammarlens}

\subsection*{Your turn}
\begin{itemize}
\raggedright
  \item \emph{Say it:} \emph{sthalam}
  \item \emph{Say it:} \emph{sthalam}, once more
  \item \emph{Say it:} \emph{naukā sthalaṃ prāpnoti.} — The boat reaches dry land.
  \item \emph{Recall:} say the Sanskrit for foam, then the Sanskrit for a bubble, then say \emph{sthalam} again
\end{itemize}

\begin{tcolorbox}[breakable,colback=teal!4,colframe=teal!35!black,title={Before you move on}]
What does \sk{स्थलम्} mean? (\textquotedblleft{}Dry land\textquotedblright{}.) Say it once more.
\end{tcolorbox}
\section[\texorpdfstring{\sk{वेला} (velā)}{वेला (velā)}]{\sk{वेला} (velā) — the coast; the shoreline}

\label{lesson:SA-C417-vela}



\begin{quote}
Before the new one: say the Sanskrit for a bubble, then the Sanskrit for dry land.
\end{quote}

\subsection*{You'll want to know: \sk{वेला}}
\textbf{\sk{वेला}} — \emph{velā} — \textquotedblleft{}the coast; the shoreline\textquotedblright{}.

\begin{quote}
\sk{वेलायां} \sk{सिकता} \sk{उष्णा} \sk{अस्ति।}
\end{quote}

\begin{quote}
\emph{velāyāṃ sikatā uṣṇā asti.}
\end{quote}

\begin{quote}
The sand on the coast is hot.
\end{quote}

A feminine noun.

\begin{cousinweb}
\textbf{\sk{वेला}} (\emph{velā}) is the word MW gives for a limit or boundary, a time, and the boundary of sea and land, the coast; MW gives no derivation, so its origin is uncertain (MW).
\end{cousinweb}

\begin{grammarlens}[title={What you've built}]
One more word for getting around.
\end{grammarlens}

\subsection*{Your turn}
\begin{itemize}
\raggedright
  \item \emph{Say it:} \emph{velā}
  \item \emph{Say it:} \emph{velā}, once more
  \item \emph{Say it:} \emph{velāyāṃ sikatā uṣṇā asti.} — The sand on the coast is hot.
  \item \emph{Recall:} say the Sanskrit for a bubble, then the Sanskrit for dry land, then say \emph{velā} again
\end{itemize}

\begin{tcolorbox}[breakable,colback=teal!4,colframe=teal!35!black,title={Before you move on}]
What does \sk{वेला} mean? (\textquotedblleft{}The coast; the shoreline\textquotedblright{}.) Say it once more.
\end{tcolorbox}
\section[(dialogue)]{Looking back: from Breakdowns to Sandbanks and Shores}

\label{lesson:SA-R417-review-410-417}



\begin{quote}
Say the Sanskrit for a hole and leaks.
\end{quote}

\subsection*{Your turn}
Say these aloud:

\begin{itemize}
\raggedright
  \item a hole, leaks, loose, split open, scrapes
  \item a carriage road, moves off, a signal, a blockage, keeps away from
  \item a halt, turns round, approaches, passes by, a kos
  \item raises up, gets ready, a sack, delivers, arranges
  \item losing one's bearings, astray, finds the way, unfamiliar, a newcomer
  \item a flowing current, a whirlpool, swims across, the far bank, a causeway
  \item a boatman, an oar, a helmsman, rows, comes ashore
  \item a sandbank, foam, a bubble, dry land, the coastline
\end{itemize}

\begin{tcolorbox}[breakable,colback=teal!4,colframe=teal!35!black,title={Before you move on}]
A hole? (\textbf{\sk{छिद्रम्}}, \emph{chidram}.) Leaks? (\textbf{\sk{क्षरति}}, \emph{kṣarati}.) Loose? (\textbf{\sk{शिथिलः}}, \emph{śithilaḥ}.) Split open? (\textbf{\sk{विदीर्णः}}, \emph{vidīrṇaḥ}.) Scrapes? (\textbf{\sk{घर्षति}}, \emph{gharṣati}.) A carriage road? (\textbf{\sk{रथ्या}}, \emph{rathyā}.) Moves off? (\textbf{\sk{अपयाति}}, \emph{apayāti}.) A signal? (\textbf{\sk{संकेतः}}, \emph{saṅketaḥ}.) A blockage? (\textbf{\sk{रोधः}}, \emph{rodhaḥ}.) Keeps away from? (\textbf{\sk{वर्जयति}}, \emph{varjayati}.) A halt? (\textbf{\sk{विरामः}}, \emph{virāmaḥ}.) Turns round? (\textbf{\sk{निवर्तते}}, \emph{nivartate}.) Approaches? (\textbf{\sk{उपगच्छति}}, \emph{upagacchati}.) Passes by? (\textbf{\sk{अतिगच्छति}}, \emph{atigacchati}.) A kos? (\textbf{\sk{क्रोशः}}, \emph{krośaḥ}.) Raises up? (\textbf{\sk{उन्नयति}}, \emph{unnayati}.) Gets ready? (\textbf{\sk{सन्नह्यति}}, \emph{sannahyati}.) A sack? (\textbf{\sk{गोणी}}, \emph{goṇī}.) Delivers? (\textbf{\sk{प्रापयति}}, \emph{prāpayati}.) Arranges? (\textbf{\sk{सम्पादयति}}, \emph{sampādayati}.) Losing one's bearings? (\textbf{\sk{दिग्भ्रमः}}, \emph{digbhramaḥ}.) Astray? (\textbf{\sk{भ्रान्तः}}, \emph{bhrāntaḥ}.) Finds the way? (\textbf{\sk{प्रतिपद्यते}}, \emph{pratipadyate}.) Unfamiliar? (\textbf{\sk{अपरिचितः}}, \emph{aparicitaḥ}.) A newcomer? (\textbf{\sk{आगन्तुकः}}, \emph{āgantukaḥ}.) A flowing current? (\textbf{\sk{स्रोतः}}, \emph{srotaḥ}.) A whirlpool? (\textbf{\sk{आवर्तः}}, \emph{āvartaḥ}.) Swims across? (\textbf{\sk{सन्तरति}}, \emph{santarati}.) The far bank? (\textbf{\sk{पारम्}}, \emph{pāram}.) A causeway? (\textbf{\sk{सेतुबन्धः}}, \emph{setubandhaḥ}.) A boatman? (\textbf{\sk{नाविकः}}, \emph{nāvikaḥ}.) An oar? (\textbf{\sk{अरित्रम्}}, \emph{aritram}.) A helmsman? (\textbf{\sk{कर्णधारः}}, \emph{karṇadhāraḥ}.) Rows? (\textbf{\sk{वाहयति}}, \emph{vāhayati}.) Comes ashore? (\textbf{\sk{उत्तरति}}, \emph{uttarati}.) A sandbank? (\textbf{\sk{पुलिनम्}}, \emph{pulinam}.) Foam? (\textbf{\sk{फेनः}}, \emph{phenaḥ}.) A bubble? (\textbf{\sk{बुद्बुदः}}, \emph{budbudaḥ}.) Dry land? (\textbf{\sk{स्थलम्}}, \emph{sthalam}.) The coastline? (\textbf{\sk{वेला}}, \emph{velā}.)
\end{tcolorbox}
